% ===================== Chapter 8 =====================
\chapter{Strings and Characters}
\label{ch:strings}

Text is everywhere in programming: names, messages, file contents, user input.
In Salam a piece of text is a \code{str}, and a single character is a \code{char}.
We have used both since Chapter~1; this chapter gathers everything about them in
one place how to write them, join them, search them, and transform them with
the standard library.

\section{String literals}

A \textbf{string literal} is text between double quotes:

\begin{salam}
func main:
    greeting str = "Hello, Salam!"
    empty str = ""
    println greeting
    println "length:", len(greeting)
end
\end{salam}

\begin{progout}
Hello, Salam!
length: 13
\end{progout}

The built-in \code{len} function reports the number of bytes in a string ---
\code{13} for \code{"Hello, Salam!"}. A string with nothing between the quotes,
\code{""}, is the \emph{empty string}; it is perfectly valid and has length 0.

\section{Strings are immutable}

A \code{str} cannot be changed in place. Every operation that seems to ``modify''
a string actually builds a \emph{new} one and leaves the original untouched. This
is why you can share strings freely without fear that someone will alter one
behind your back, and it is why joining strings always produces a fresh result
rather than growing an existing one.

\section{Joining strings with \texttt{+}}

The \code{+} operator concatenates glues together two strings, producing a
new string:

\begin{salam}
func main:
    first str = "Salam"
    last str = "World"
    full auto = first + ", " + last + "!"
    println full
end
\end{salam}

\begin{progout}
Salam, World!
\end{progout}

Each \code{+} builds a larger string from its two operands. Because the result is
itself a string, you can chain as many as you like, as the example does.

\begin{warn}
\code{+} joins a string to a string. It does \emph{not} automatically turn a
number into text: \code{"count: " + 5} is a type error, because \code{5} is an
\code{i32}, not a \code{str}. To attach a number, convert it first the
standard library's \code{str.FromInt} does exactly that, as we will see shortly.
(When you simply want to \emph{print} mixed values, \code{println} already
accepts a comma-separated mix of any types no conversion needed.)
\end{warn}

\section{Escape sequences}

Some characters are awkward to type directly inside a string a newline, a tab,
or a double quote (which would otherwise end the string). For these you use an
\textbf{escape sequence}: a backslash followed by a code.

\begin{center}
\begin{tabular}{@{}cl@{}}
\toprule
\textbf{Escape} & \textbf{Means} \\
\midrule
\code{\textbackslash n} & newline (move to the next line) \\
\code{\textbackslash t} & tab \\
\code{\textbackslash\textbackslash} & a literal backslash \\
\code{\textbackslash"} & a literal double quote \\
\bottomrule
\end{tabular}
\end{center}

\begin{salam}
func main:
    println "line one\nline two"
    println "name:\tSalam"
    println "she said \"hi\""
    println "a backslash: \\"
end
\end{salam}

\begin{progout}
line one
line two
name:	Salam
she said "hi"
a backslash: \
\end{progout}

The \code{\textbackslash n} split the first string across two lines; the
\code{\textbackslash t} inserted a tab; \code{\textbackslash"} put real quotes
\emph{inside} the text; and \code{\textbackslash\textbackslash} produced a single
backslash.

\section{Triple-quoted strings}

When a string itself contains many double quotes, escaping each one is tedious. A
\textbf{triple-quoted} string, wrapped in \code{"""}, lets you write double
quotes directly without escaping:

\begin{salam}
func main:
    line str = """She said "hi" to me"""
    println line
end
\end{salam}

\begin{progout}
She said "hi" to me
\end{progout}

Triple-quoted strings are handy for text with embedded quotes snippets of
other code, JSON, or dialogue.

\section{Characters}

A \code{char} is a single character, written between \emph{single} quotes. As we
saw in Chapter~2, a \code{char} is really a small number (its byte code), and you
can convert between the two with \code{as}:

\begin{salam}
func main:
    letter char = 'S'
    println letter              // S
    println letter as i32       // 83, the code for 'S'
    println 33 as char          // '!', the character with code 33
end
\end{salam}

\begin{progout}
S
83
!
\end{progout}

\begin{note}
\code{'S'} (single quotes) is a \code{char} one character. \code{"S"} (double
quotes) is a \code{str} text that happens to be one character long. They are
different types and are not interchangeable; keep an eye on which quotes you are
typing.
\end{note}

\section{Strings are not indexed with brackets}

You might expect \code{s[0]} to give the first character of a string, as it would
for an array. In Salam it does \emph{not} a \code{str} is not an array, and
indexing one is an error. To work with parts of a string, you use the functions
of the \code{str} standard-library module, which we turn to now.

\section{The \texttt{str} module}

The built-in \code{+} and \code{len} cover the basics, but most real string work
goes through the \code{str} module. Import it with \code{import "str"} and call
its functions with the \code{str.} prefix. Here are the ones you will use most.

\subsection{Case and trimming}

\begin{salam}
import "str"

func main:
    println str.Upper("salam")        // SALAM
    println str.Lower("SALAM")        // salam
    println str.Trim("   hi   ")      // "hi" (surrounding spaces removed)
    println str.Reverse("Salam")      // malaS
    println str.Repeat("ab", 3)       // ababab
end
\end{salam}

\begin{progout}
SALAM
salam
hi
malaS
ababab
\end{progout}

\subsection{Searching and testing}

\begin{salam}
import "str"

func main:
    println str.Find("hello world", "world")       // 6 (the index)
    println str.Contains("hello", "ell")           // true
    println str.StartsWith("hello.salam", "hello") // true
    println str.EndsWith("hello.salam", ".salam")  // true
end
\end{salam}

\begin{progout}
6
true
true
true
\end{progout}

\code{str.Find} returns the index at which the substring begins (here \code{6},
since \code{"world"} starts at position 6 in \code{"hello world"}), or a negative
value if it is not present.

\subsection{Substrings and replacement}

\begin{salam}
import "str"

func main:
    println str.Substr("programming", 0, 7)    // "program"
    println str.Replace("a-b-c", "-", "+")     // "a+b+c"
end
\end{salam}

\begin{progout}
program
a+b+c
\end{progout}

\code{str.Substr(s, start, n)} extracts \code{n} characters beginning at index
\code{start} so \code{Substr("programming", 0, 7)} takes seven characters from
the start. This is also how you read a single character: \code{str.Substr(s, i,
1)} gives the one-character string at index \code{i}.

\subsection{Converting between strings and numbers}

Turning numbers into text and back is constant work, and the \code{str} module
makes it direct:

\begin{salam}
import "str"

func main:
    n i32 = 2026
    text auto = str.FromInt(n as i64)     // "2026"
    println "Year: " + text               // now '+' works: both are strings

    back auto = str.ToInt("42")           // 42, an i32
    println back + 8                      // 50
end
\end{salam}

\begin{progout}
Year: 2026
50
\end{progout}

\code{str.FromInt} converts an integer to its text form (it takes an \code{i64},
so cast smaller integers with \code{as i64}); \code{str.ToInt} parses a string
back into an \code{i32}. With \code{FromInt} you can finally use \code{+} to
attach a number to a message, because both sides are now strings.

\subsection{Splitting and joining}

Two of the most useful string operations break a string apart and put one back
together. \code{str.Split} cuts a string at every occurrence of a separator,
returning a \code{Vector<str>}; \code{str.Join} does the reverse:

\begin{salam}
import "str"

func main:
    mut parts Vector<str> = str.Split("ada,grace,linus", ",")
    println "count:", parts.len()
    println str.Join(parts, " & ")
    parts.free()
end
\end{salam}

\begin{progout}
count: 3
ada & grace & linus
\end{progout}

Splitting \code{"ada,grace,linus"} on \code{","} yields a vector of three
strings; joining them with \code{" \& "} stitches them back into one. Because
\code{Split} returns a vector, remember to \code{free} it when you are done
(Chapter~7).

\section{A worked example: counting words}

Putting several of these together, here is a small program that counts the words
in a sentence and prints each on its own line:

\begin{salam}
import "str"

func main:
    sentence str = "the quick brown fox"
    mut words Vector<str> = str.Split(sentence, " ")
    println "word count:", words.len()
    mut i i32 = 0
    while i < words.len():
        println str.Upper(words.get(i)[0])
        i += 1
    end
    words.free()
end
\end{salam}

\begin{progout}
word count: 4
THE
QUICK
BROWN
FOX
\end{progout}

\begin{deep}[In Depth: bytes, characters, and Unicode]
A Salam \code{str} is stored as a sequence of bytes in UTF-8 encoding, and
\code{len} reports the number of \emph{bytes}. For plain English text, one byte
is one character, so the distinction never bites. But a character outside the
basic Latin range an accented letter, or a letter of another script may
occupy several bytes, so its byte length and its visible length differ. The
\code{str} module's higher-level functions (\code{Split}, \code{Find},
\code{Replace}, and friends) are the safe way to manipulate text; reach for raw
byte arithmetic only when you know your data is plain ASCII.
\end{deep}

\section{Chapter summary}

\begin{itemize}[leftmargin=1.4em]
  \item A \code{str} is immutable text in double quotes; \code{len(s)} gives its
        byte length; \code{+} concatenates.
  \item Escape sequences (\code{\textbackslash n}, \code{\textbackslash t},
        \code{\textbackslash"}, \code{\textbackslash\textbackslash}) embed
        special characters; \code{"""..."""} avoids escaping inner quotes.
  \item A \code{char} is a single character in single quotes and converts to/from
        its numeric code with \code{as}.
  \item Strings are \emph{not} indexed with \code{[]}; use the \code{str} module.
  \item \code{str} provides \code{Upper}, \code{Lower}, \code{Trim},
        \code{Reverse}, \code{Repeat}, \code{Find}, \code{Contains},
        \code{StartsWith}, \code{EndsWith}, \code{Substr}, \code{Replace},
        \code{FromInt}, \code{ToInt}, \code{Split}, and \code{Join}.
\end{itemize}

\section{Exercises}

\exercise Ask for a \code{str} variable holding your full name and print its
length, its uppercase form, and its reverse.

\exercise Build the string \code{"Item No.7 costs \$42"} by converting the numbers
7 and 42 with \code{str.FromInt} and joining with \code{+}.

\exercise Given \code{"one,two,three,four"}, split on the comma, print how many
parts there are, and print them re-joined with \code{" - "}.

\exercise Write a function \code{is\_palindrome(s str) bool} that returns whether
a string reads the same forwards and backwards (hint: compare \code{s} with
\code{str.Reverse(s)}).

\exercise Use \code{str.Substr} to print each character of \code{"Salam"} on its
own line.

\exercise Write a program that, given a filename like \code{"report.txt"}, prints
just the extension (\code{"txt"}) by finding the last \code{"."}.
